% \vspace{-2mm}
\section{Conclusions}
\label{sec:conclusions}

% In this study, we %have focused 
% focus on the long distance-aware multi-hop reasoning task that has rarely been %ignored 
% approached in the previous studies. %We have proposed a 
% Our proposed %novel 
% model %, called 
% GMH, which seamlessly unifies both local path factors and global embedding factors, %, to approach this task which yields considerable performance improvements on both short and long distance scenarios.
% can successfully accomplish this task as demonstrated by extensive experiments.
% %the newly constructed dataset.
% %Our future work will consider incorporating external text data and developing our model to solve further complex tasks.
We have studied the multi-hop reasoning task in long distance scenarios and proposed a general model which could tackle both short and long distance reasoning scenarios.
Extensive experiments showed the effectiveness of our model on three benchmarks.
% We hope insights from this work will inspire future developments of KG completion and KG-related applications.
We will further consider the feasibility of applying our model to complex real-world datasets with more long distance reasoning scenarios and more relation types.
% The superior performance of our model in long distance reasoning scenarios provides opportunities for conducting complex reasoning on KG-related applications, e.g., complex question answering based on KG.
Besides, we have noticed that there are other ``interference'' in long distance reasoning. For example, noise from the KG itself, i.e., the fact that it lacks validity. These noises can gradually accumulate during long distance reasoning and affect the result confidence. We leave the further investigation to future work.


\section{Acknowledgements} 
    We sincerely thank Jun Wang and Xu Zhang for their constructive suggestions on this paper.
    This work was supported by the China Postdoctoral Science Foundation (No.2021TQ0222).